\section{The class-field functor and the actual adelic coefficient sheaf over the split base}
\label{sec:connes-support-functor}

We use Alain Connes and Caterina Consani's original author source
\emph{Knots, primes and class field theory}, arXiv:2501.06560v1
\cite{ConnesConsani2025}. The inputs are their associated covers in
\texttt{sec2}, the class-field functor \texttt{main}, and, in Section 4,
\texttt{localv}, \texttt{structure}, and \texttt{structure1}.
The source archive and its author TeX are retained without alteration.
The following constructions apply the full spectrum (FL3) to those
objects. In particular they do not replace the support spectrum by
one selected branch. All sheaf cohomology below is cohomology of the
underlying complex vector sheaf; a multiplication is asserted only
where its compatibility is proved.

\subsection{Extending the cover functor, with the inertia retained}

Put $S=\SpecLat L$, where $L$ is a finite nontrivial bounded
distributive lattice. For a topological space $Z$ define
\[
 T_SZ=S\sqcup Z,\qquad
 \operatorname{Open}(T_SZ)=\operatorname{Open}(S)
 \cup\{S\cup U: \varnothing\ne U\in\operatorname{Open}(Z)\}.
 \tag{CCF1}
\]
This is a topology: a union involving an old open contains all of
$S$, and an intersection of two such opens either has the displayed
form or equals $S$. It is functorial in continuous maps on both
strata. The inverse image of $S\cup U$ under $f\sqcup h$ is
$S'\cup h^{-1}U$, interpreted as the open $S'$ when $h^{-1}U$ is
empty. The foundational prime classification gives the actual
homeomorphism $\Spec G_L(R)=T_S\Spec R$ for every nonzero ring $R$.

Use the original spaces and group
\[
 Y_{\Q}=\Q^\times\backslash\mathbb A_{\Q},\quad
 K=\widehat{\mathbb Z}^{\times},\quad X_{\Q}=Y_{\Q}/K,
 \quad \chi:K\twoheadrightarrow G,
 \tag{CCF2}
\]
where $G$ is finite abelian and $\chi$ is continuous. Connes--Consani's
cover is $X_{\Q}^{\chi}=(Y_{\Q}\times G)/K$, with
$w\cdot(y,g)=(yw,\chi(w)g)$. Extend the $K$ action trivially on $S$.
There are canonical, $G$-equivariant homeomorphisms
\[
 (T_SY_{\Q})/K\simeq T_SX_{\Q},\qquad
 ((T_SY_{\Q})\times G)/K\simeq T_SX_{\Q}^{\chi}.
 \tag{CCF3}
\]
For the first statement, a saturated old open is the inverse image
of an old quotient open, and a support open is unchanged. For the
second, write an open in $(T_SY_{\Q})\times G$ by its $G$ slices.
Invariance under $K$ and surjectivity of $\chi$ imply that either
all slices have no old point, with the same support open, or every
slice contains all of $S$. In the latter case their old parts form
an open invariant subset of $Y_{\Q}\times G$. This proves the
quotient topology in (CCF3), not only a bijection of points.

The map in (CCF3) over an old point has the original fibre and
stabilizer. In particular, for a representative $a$ with
$a_\infty\ne0$, put $Z(a)=\{v:a_v=0\}$. Its fibre and inertia are
\[
 G/\chi(H_{Z(a)}),\qquad \chi(H_{Z(a)}),\qquad
 H_{Z(a)}=\prod_{v\in Z(a)}\mathbb Z_v^\times\subset K.
 \tag{CCF4}
\]
Indeed $aw=qa$ forces $q=1$ at the real coordinate; at a nonzero
finite coordinate it forces $w_v=1$, and at a zero coordinate it
imposes no condition. This proves the stabilizer formula directly.
For a support point $s\in S$ the fibre is one coarse point and its
$G$ stabilizer is all of $G$. Thus the extension changes the
inertia geometry along every support point.

The group action retained at that point is explicit. The action
groupoid of $K$ on $\{s\}\times G$ is equivalent to the one-object
groupoid $B\ker\chi$: the orbit is transitive and the automorphisms
of $(s,1)$ are exactly $\ker\chi$. Its map to the base groupoid
$BK$ is induced by $\ker\chi\hookrightarrow K$. Consequently
the finite covering data at the new point are
\[
 B\ker\chi\longrightarrow BK,\qquad K/\ker\chi\simeq G.
 \tag{CCF5}
\]
This also computes the exact transfer on constant complex
coefficients. The induced coefficient space is $\mathbb C[G]$;
the unit sends $z$ to the constant vector $(z)_{g\in G}$, and the
counit sums its entries. Their composite is $|G|z$. Coarse fibre
cardinality one therefore does not compute this transfer: the
stabilizer data in (CCF5) give its degree $|G|$.

For a finite abelian extension $E/\Q$, take the source's reciprocity
map $\chi_E:K\twoheadrightarrow G_E$. Given an embedding
$\sigma:E_1\to E_2$, write it as inclusion followed by $h\in G_{E_1}$,
with restriction $r:G_{E_2}\to G_{E_1}$. The source's map is
\[
 \rho_\sigma(g)=r(g)h,\qquad
 [y,g]\longmapsto[y,\rho_\sigma(g)].
 \tag{CCF6}
\]
It descends since $r\chi_{E_2}=\chi_{E_1}$ and both groups are
abelian. Extend it by the identity on $S$. The slice proof above
shows continuity, and composition agrees with composition of the
$\rho_\sigma$. Hence (CCF3) defines an actual contravariant
extension of the supplied class-field functor. The reciprocity
theorem used to supply $\chi_E$ is the human-source input; no
new proof of class field theory is asserted here.

There is a compatible algebraic action on the original base.
Every $g\in G_E$ acts on $G_L(\mathcal O_E)$ by
$(a,\lambda)\mapsto(g(a),\lambda)$. On prime spectra it fixes
each support prime and acts on the arithmetic primes by the usual
Galois action. Its map to $\Spec G_L(\mathbb Z)$ is the identity
on $S$ and the original arithmetic prime map on the other stratum.
This follows by taking inverse images in (FL3). In particular the
fixed support points in (CCF3) have the same coefficient-action
origin as the fixed support primes of the actual semiring.

The original periodic orbit $C_p$ and its mapping torus remain
old subspaces, with their original topology. In the source's
orientation their period and monodromy are
\[
 \operatorname{length}(C_p)=\log p,\qquad g\longmapsto\chi_E(p)g
 \quad\text{when }\chi_E(\mathbb Z_p^\times)=1.
 \tag{CCF7}
\]
This is also obtained directly from
$(h,\lambda)\sim(ph,p\lambda)$ in the source mapping torus.
The new support stratum is fixed by the scaling action. The
unramified monodromy in (CCF7) is a permutation on $\mathbb C[G_E]$:
its eigenvalues are roots of unity. The new stabilizers, transfers,
and the cohomology calculated next must all be retained in applying
this observation to the split base.

\subsection{A complete resolution of the original semilocal sheaves}

Let $X=\Spec\mathbb Z$, with generic point $\eta$. In this subsection
$F$ denotes a finite set of finite primes, not the support space $S$.
Retain the actual Bruhat--Schwartz functions and group algebras
\[
 V=\mathcal S(\mathbb A_{\Q}),\qquad
 B=V\rtimes\mathbb Q^\times,\qquad
 \mathcal O(X\setminus F)=\mathcal S(\mathbb A_{F\cup\{\infty\}}),
 \quad
 \mathcal A(X\setminus F)=\mathcal O(X\setminus F)
               \rtimes\mathbb Z[F^{-1}]^\times.
 \tag{CCF8}
\]
Every inclusion into the generic space inserts $1_{\mathbb Z_p}$
at the missing finite primes. These are exactly the sheaves
\texttt{structure} and \texttt{structure1} of Connes--Consani,
including their specified transition maps.

Write $V^{(p)}$ for the restricted tensor product over all places
other than $p$, including the real factor, and put
\[
 V_p=\mathbb C1_{\mathbb Z_p}\otimes V^{(p)},\qquad
 B_p=\bigoplus_{r\in\mathbb Q^\times:\ v_p(r)=0}V_pU_r.
 \tag{CCF9}
\]
They are the actual stalks of $\mathcal O$ and $\mathcal A$ at
the closed point $p$. To see this, the neighbourhoods of $p$ may
exclude any finite set not containing $p$. Their union allows
arbitrary finite coefficients away from $p$ but retains precisely
the displayed unramified factor and the condition $v_p(r)=0$.
Define the marked local defect spaces
\[
 D_V=\bigoplus_p V/V_p,\qquad D_B=\bigoplus_p B/B_p,
 \quad \delta_V(f)=([f]_p)_p,\quad\delta_B(b)=([b]_p)_p.
 \tag{CCF10}
\]
The maps land in direct sums: a Bruhat--Schwartz function involves
only finitely many nonstandard finite factors, and a rational
number has nonzero valuation at only finitely many primes.

There are exact flasque resolutions
\[
 \begin{aligned}
 0&\longrightarrow\mathcal O\longrightarrow c_XV
       \longrightarrow\bigoplus_p i_{p*}(V/V_p)\longrightarrow0,\\
 0&\longrightarrow\mathcal A\longrightarrow c_XB
       \longrightarrow\bigoplus_p i_{p*}(B/B_p)\longrightarrow0.
 \end{aligned}
 \tag{CCF11}
\]
The displayed arrows are the quotient maps at every included
closed point. Their kernel on $X\setminus F$ is exactly the
space in (CCF8): a coefficient is in every retained $V_p$ exactly
when it has the unramified factor at every $p\notin F$, and a
group label must be a unit at those same primes. The source's
local characterization \texttt{localv}, or its factorization
$f(x,y)=1_{\mathbb Z_p}(x)f(0,y)$, proves this coefficient claim.
At $p$ the cokernel arrow is onto by definition, and at $\eta$
its target stalk is zero. This proves sheaf exactness.

For clarity, the direct sums in (CCF11) have finite-support
sections on each nonempty open. A section has zero generic germ,
so vanishes on a smaller cofinite open. The remaining support is
finite. Conversely any finite collection of stalk values defines
a section of the corresponding sum of skyscraper sheaves. Its
restriction map forgets the excluded components and is onto.
These sheaves are therefore flasque. The constant sheaves are
flasque because every nonempty open of $X$ is connected and
their nonempty restrictions are identities. It follows that
\[
 R\Gamma(X,\mathcal O)=[V\xrightarrow{\delta_V}D_V],\qquad
 R\Gamma(X,\mathcal A)=[B\xrightarrow{\delta_B}D_B],
 \tag{CCF12}
\]
in degrees zero and one, with all higher groups zero. In
particular their degree-zero groups are exactly
$\mathcal S(\mathbb R)$ and $\mathcal S(\mathbb R)\rtimes\{\pm1\}$,
as in the original source. This calculation also determines the
previously unspecified higher sheaf group.

Here is its explicit mixed-place content. Define
$E_p=\ker(f\mapsto f(0):\mathcal S(\mathbb Q_p)\to\mathbb C)$.
Since $1_{\mathbb Z_p}(0)=1$, the original space splits as
$\mathbb C1_{\mathbb Z_p}\oplus E_p$. Expanding each finite
tensor therefore gives, with all factors retained,
\[
 V=\bigoplus_{T\text{ finite}}V_T,\qquad
 V_T=\mathcal S(\mathbb R)\otimes\bigotimes_{p\in T}E_p.
 \tag{CCF13}
\]
Factors outside $T$ are $1_{\mathbb Z_p}$. For $r\in\mathbb Q^\times$
put $D(r)=\{p:v_p(r)\ne0\}$ and $T_r=T\cup D(r)$. Then (CCF12)
is the direct sum of the actual diagonal complexes
\[
 [V_T\longrightarrow V_T^{\,T}]
 \quad\text{or}\quad
 [V_TU_r\longrightarrow (V_TU_r)^{\,T_r}],
 \tag{CCF14}
\]
respectively. An empty target is zero. Indeed $V_TU_r$ survives
in $B/B_p$ precisely when $p\in T_r$; the quotient sends its
coefficient unchanged to that marked component. Thus
\[
 \begin{aligned}
 H^1(X,\mathcal O)&=\bigoplus_{|T|\ge2}
       V_T\otimes(\mathbb C^T/\Delta\mathbb C),\\
 H^1(X,\mathcal A)&=\bigoplus_{r,T:\ |T_r|\ge2}
       V_TU_r\otimes(\mathbb C^{T_r}/\Delta\mathbb C).
 \end{aligned}
 \tag{CCF15}
\]
Equations (CCF11)--(CCF15) retain every mixed finite-place
coefficient and every original rational label. The decomposition
(CCF13) is not claimed to be invariant under nonunit finite
scalings; their exact defect is computed below.

\subsection{Relative cohomology after adjoining the whole support spectrum}

For either $(\mathcal H,W,D)=(\mathcal O,V,D_V)$ or
$(\mathcal A,B,D_B)$ take the universal extension $U_S\mathcal H$
of (FL12). Its restriction to $S$ is the constant sheaf $c_SW$,
and its arithmetic-to-support map is the actual generic
inclusion, followed by the constant diagonal. If $S$ has one
point, (CCF12) and the relative cone give the particularly
simple complete answer
\[
 R\Gamma_X(T_SX,U_S\mathcal H)\simeq D[-1].
 \tag{CCF16}
\]
Here $D[-1]$ means $D$ in cohomological degree one. To prove the
identity with signs specified, use the complex
$W\to D\oplus W$, $w\mapsto(\delta w,w)$.
The map $(d,w)\mapsto d-\delta w$ kills this differential and
induces an isomorphism on its cokernel. Its kernel consists
exactly of $(\delta w,w)$; its degree-zero kernel is zero.
This is a full quasi-isomorphism, not a subtraction of dimensions.

For arbitrary $S$, let $C^\bullet(S,W)$ be the complete finite
cochain complex (FL18), with its stated specialization orientation.
An explicit relative complex is
\[
 \begin{aligned}
 K^0&=W,& K^1&=D\oplus C^0(S,W),&
 K^{n}&=C^{n-1}(S,W)\quad(n\ge2),\\
 d^0w&=(\delta w,\Delta w),& d^1(d,c)&=d_Sc,&
 d^{n}&=d_S\quad(n\ge2).
 \end{aligned}
 \tag{CCF17}
\]
Changing the support coordinate signs in the mapping fibre
gives precisely this convention. Since $d_S\Delta=0$, it is
a complex. Let $\widetilde C^\bullet(S,W)$ denote
$C^\bullet(S,W)$ modulo the constant diagonal in degree zero.
Projection in (CCF17) gives an exact sequence of complexes
whose kernel is $[W\to D\oplus\Delta W]$. The preceding explicit
contraction identifies that kernel with $D[-1]$. Therefore
\[
 \begin{gathered}
 H_X^0=0,\qquad
 0\longrightarrow D\longrightarrow H_X^1
   \longrightarrow W\otimes\widetilde H^0(S,\mathbb C)
   \longrightarrow0,\\
 H_X^{j+1}=W\otimes H^j(S,\mathbb C)\quad(j\ge1).
 \end{gathered}
 \tag{CCF18}
\]
No support component or support cycle has been suppressed.
Choosing $s_0\in S$ splits the degree-one sequence by the
retraction $[(d,c)]\mapsto d-\delta(c_{s_0})$. It is well defined
on closed representatives modulo $(\delta w,\Delta w)$.
The choice is stated explicitly; the injection of $D$ is canonical.

The original arithmetic map is also explicit. Forgetting the
support part in (CCF17) sends $[(d,c)]$ to $[d]$ in
$D/\delta W=H^1(X,\mathcal H)$. On the canonical $D$ summand it
is the quotient map
\[
 D\longrightarrow D/\delta W.
 \tag{CCF19}
\]
It is onto. Thus the image in the original semilocal sheaf
cohomology is exactly (CCF15). In particular the vanishing image
computed for the different coefficient sheaf $\mathcal O_X$ in
(AQ17) cannot be substituted here: (CCF19) computes the required
comparison for these actual Connes--Consani coefficients.

\subsection{Both pole characters and Fourier reflection on the new groups}

Use additive Haar measure with $\operatorname{vol}(\mathbb Z_p)=1$
at finite places and $dx$ at the real place. Define on $V$
\[
 \epsilon_0(f)=f(0),\qquad \epsilon_1(f)=\int_{\mathbb A_\Q}f(x)\,dx,
 \qquad (\vartheta_\lambda f)(x_\infty,x_f)
       =f(\lambda^{-1}x_\infty,x_f),\quad\lambda>0.
 \tag{CCF20}
\]
Insertion of $1_{\mathbb Z_p}$ preserves both functionals, so
they define a sheaf morphism
$\mathcal O\to c_X\mathbb C^2$. It is onto on every nonempty
open. For example $e^{-\pi x^2}$ has moments $(1,1)$ and
$x^2e^{-\pi x^2}$ has moments $(0,1/(2\pi))$, and both have
standard factors at every finite place. The usual Gaussian
integral gives these values; it also follows by differentiating
$\int e^{-a\pi x^2}dx=a^{-1/2}$ at $a=1$.
The exact pole action is
\[
 (\epsilon_0,\epsilon_1)\vartheta_\lambda
       =\operatorname{diag}(1,\lambda)(\epsilon_0,\epsilon_1).
 \tag{CCF21}
\]
The first equality is evaluation and the second is change of
variables; no factor has been inferred from the height of a prime.

Put $\mathcal O^0=\ker(\epsilon_0,\epsilon_1)$ and
$V^0=\ker((\epsilon_0,\epsilon_1):V\to\mathbb C^2)$.
At every finite prime the canonical map
\[
 V^0/(V^0\cap V_p)\xrightarrow{\ \sim\ }V/V_p
 \tag{CCF22}
\]
is an isomorphism. Injectivity is its displayed kernel. For
surjectivity, subtract from a given $f$ an element of $V_p$
with the same two moments, using the two real test functions
above. This calculation retains the two pole conditions rather
than dropping them without a map. For one support point,
$\mathcal O^0$ and $\mathcal O$ therefore have the same relative
degree-one group $D_V$. For general $S$, their difference is
exactly the relative complex of the constant pole sheaf:
\[
 R\Gamma_X(U_Sc_X\mathbb C^2)
       =\widetilde C^\bullet(S,\mathbb C^2)[-1],
 \tag{CCF23}
\]
with the degree convention of (CCF17). This follows either
by the same contraction with $D=0$, or by the short exact
sequence of the three coefficient sheaves. Its two scaling
characters are still $1$ and $\lambda$ on every support term.

Take the original additive Fourier transform, with character
$e^{-2\pi ixy}$ at infinity and the compatible conductor
$\mathbb Z_p$ characters at finite places. Character
orthogonality gives $\widehat{1_{\mathbb Z_p}}=1_{\mathbb Z_p}$:
the integral over $\mathbb Z_p$ is one when the character is
trivial there and zero otherwise, by translating it by an
element with nontrivial character value. The tensor Fourier
transform $\mathcal F$ thus preserves every $V_p$, every
semilocal transition, and $V^0$. It exchanges the two moments
and obeys the exact identities
\[
 \mathcal F^2f(x)=f(-x),\qquad
 \mathcal F\vartheta_\lambda
     =\lambda\vartheta_{\lambda^{-1}}\mathcal F.
 \tag{CCF24}
\]
The second follows by $x_\infty=\lambda t$ in the defining
integral. The first is Fourier inversion on the real Schwartz
space and on finite local test functions; the tensor identity
then holds on every original adelic test. Consequently
(CCF24) holds on $D_V$, all complexes (CCF17), and their
arithmetic images (CCF19). It is an actual weight-one reflection
on the reconstructed objects, including their support cycles.

There is a precise multiplicative target for the crossed
product as well. Fourier transforms pointwise multiplication
to additive convolution. If $\alpha_r f(x)=f(r^{-1}x)$, then
the product formula gives $|r|_{\mathbb A}=1$, and
$\mathcal F\alpha_r=\alpha_{r^{-1}}\mathcal F$. Hence
\[
 \sum f_rU_r\longmapsto\sum\widehat f_r\,U_{r^{-1}}
 \tag{CCF25}
\]
is an isomorphism from the pointwise crossed-product sheaf
to the convolution crossed-product sheaf. In fact it sends
$f_r\alpha_r(g_s)U_{rs}$ to
$\widehat f_r*\alpha_{r^{-1}}(\widehat g_s)U_{r^{-1}s^{-1}}$,
which is exactly the target product. At a semilocal stage the
same product formula holds for its unit group, and Fourier
fixes the inserted standard factors. Thus (CCF25) respects
every original restriction and induces the corresponding
relative-cohomology isomorphism. Its scaling relation is still
(CCF24), with the transported convolution action
$\lambda\vartheta_{\lambda^{-1}}$ written in full.

\subsection{The exact positive receiver and the nonunit-prime defect}

The additive Haar metric gives a concrete positive receiver
on (CCF16). It is not introduced by assigning a weight to a
topological degree. In the $p$ factor, the orthogonal projection
onto the original unramified line is
\[
 (P_pf)(x_p,x^{(p)})=1_{\mathbb Z_p}(x_p)
                \int_{\mathbb Z_p}f(t,x^{(p)})\,dt.
 \tag{CCF26}
\]
It preserves Bruhat--Schwartz functions and has image $V_p$.
Consequently $[f]_p\mapsto(1-P_p)f$ is a well-defined
isomorphism $V/V_p\to(1-P_p)V$. Give this actual quotient
the positive norm $\|(1-P_p)f\|_{L^2(\mathbb A)}$ and give
$D_V$ the orthogonal direct-sum norm. It vanishes only on the
zero quotient class, since $(1-P_p)f=0$ implies $f\in V_p$.

On $D_B$, apply the same quotient projection to a coefficient
$f_r$ when $v_p(r)=0$, and leave it unchanged when
$v_p(r)\ne0$, then sum squared norms over $(p,r)$.
This is exactly the quotient by the subspace $B_p$ of (CCF9),
with counting measure on rational group labels. All sums are
finite on the algebraic spaces under discussion. Both forms
satisfy
\[
 \langle\vartheta_\lambda d,\vartheta_\lambda d'\rangle
       =\lambda\langle d,d'\rangle,
 \tag{CCF27}
\]
because $P_p$ acts only on a finite-place coordinate and real
Haar measure changes by $\lambda$. Fourier is an isometry
for these forms and commutes with $P_p$; (CCF25) also permutes
the rational labels without changing their unit condition.
For the full support space, any chosen $s_0$ in (CCF18),
together with a fixed positive form on the finite-dimensional
support cohomology, extends (CCF27) to every relative group.
The choice of the support form and splitting is part of that
receiver and is not asserted to be canonical.

The receiving map (CCF19) has its actual kernel $\delta W$.
Its attained positive quotient is calculated using the closure
of this kernel. Let $\overline D_V$ be the Hilbert completion
of the above orthogonal sum. If $F$ contains all nonstandard
finite factors of a test function $f$, the commuting projections
give the exact orthogonal decomposition
\[
 f=\sum_{T\subset F}
       \left(\prod_{p\in T}(1-P_p)\prod_{p\in F\setminus T}P_p\right)f.
 \tag{CCF28}
\]
This follows by expanding $\prod_{p\in F}(P_p+(1-P_p))=I$;
different choices are orthogonal in at least one factor.
For two different primes and any test $f$, set
$R_{pq}=(1-P_p)(1-P_q)f$. It is the same mixed function
whether extracted from the $p$ or $q$ component of $\delta_Vf$.
Thus the bounded maps on $\overline D_V$
\[
 d\longmapsto(1-P_q)d_p-(1-P_p)d_q
 \tag{CCF29}
\]
vanish on $\delta_VV$ and on its closure. They may be nonzero:
take $d_p=h_p\otimes h_q\otimes h_\infty$, $d_q=0$, with
$h_p\perp1_{\mathbb Z_p}$, $h_q\perp1_{\mathbb Z_q}$ and
nonzero Schwartz factors. This also proves directly that the
arithmetic mixed image in (CCF19) is not lost under completion.

An exact orthogonal decomposition computes the whole quotient.
Put $E_p^{\perp}=\ker(\int_{\mathbb Z_p})$ in the local
Bruhat--Schwartz space. Replace $E_p$ in (CCF13) by this
orthogonal complement, and call the resulting summand
$V_T^{\perp}$. All unramified factors have norm one. Then
\[
 \overline D_V=\widehat{\bigoplus}_{T\ne\varnothing}
                  \overline V_T^{\perp}\otimes\mathbb C^T,
 \quad
 \overline{\delta_VV}=\widehat{\bigoplus}_{T\ne\varnothing}
                  \overline V_T^{\perp}\otimes\Delta\mathbb C,
 \tag{CCF30}
\]
and the attained Hilbert quotient is
\[
 \overline D_V/\overline{\delta_VV}
   \simeq\widehat{\bigoplus}_{|T|\ge2}
      \overline V_T^{\perp}\otimes(\Delta\mathbb C)^{\perp}.
 \tag{CCF31}
\]
To prove these identities, expand any finite tensor in the
mutually orthogonal choices $\mathbb C1_{\mathbb Z_p}$ and
$E_p^\perp$. Its $T$ component appears in $D_V$ exactly at
the marked primes $p\in T$, and $\delta_V$ is the diagonal
on that component. Such finite tensors are dense in the
restricted adelic $L^2$ tensor product, and finite collections
of their diagonal images are obtained from one finite sum
in $V$. Taking closures therefore gives exactly (CCF30).
Orthogonal projection onto the diagonal in each finite
$\mathbb C^T$ gives (CCF31). The projection is the literal
arithmetic mean $|T|^{-1}\sum_{p\in T}$, so all multiplicities
are retained. The natural algebraic map from (CCF15) into
(CCF31) is injective: a finite component has zero projection
precisely when its marked values are diagonal, already the
algebraic relation. This is the complete quotient calculation,
not a supposition that the image of a dense subspace is closed.

Exactly the same proof for $B$ replaces $T$ by $T_r=T\cup D(r)$
and adds its original label $r$. It yields
\[
 \overline D_B/\overline{\delta_BB}
   \simeq\widehat{\bigoplus}_{r,T:\ |T_r|\ge2}
       \overline V_T^{\perp}U_r\otimes H_{T_r},\qquad
 H_{T_r}=\{(c_p)_{p\in T_r}:\sum_{p\in T_r}c_p=0\}.
 \tag{CCF32}
\]
The finite factor uses its inherited Euclidean metric. Equations
(CCF27) and (CCF24) pass to both attained quotients because
the diagonal relation and its closure are invariant.

Finally the operation of a finite prime must be traced through
these maps, rather than inferred from (CCF27). Let
$(\alpha_pf)(x)=f(p^{-1}x)$ act at all adelic places. The
product formula makes it an isometry on $L^2(\mathbb A)$,
but it does not preserve $V_p$: its local standard factor
becomes $1_{p\mathbb Z_p}$. The exact nonunit defect is the map
\[
 \partial_p:V_p\longrightarrow V/V_p,\qquad
      f\longmapsto[\alpha_pf]_p.
 \tag{CCF33}
\]
For $f=1_{\mathbb Z_p}\otimes g$ it is represented by
$1_{p\mathbb Z_p}\otimes\alpha_p^{(p)}g$; subtracting its
orthogonal unramified projection gives the exact vector
\[
 (1_{p\mathbb Z_p}-p^{-1}1_{\mathbb Z_p})
                  \otimes\alpha_p^{(p)}g.
 \tag{CCF34}
\]
Its squared norm is
$(p^{-1}-p^{-2})\|\alpha_p^{(p)}g\|^2
=(1-p^{-1})\|g\|^2$, because the product of the other local
moduli of the rational $p$ is $p$. Thus the defect is nonzero
and has a calculated norm on the original coordinates.
It is an actual element of the supported relative group,
not a discarded failure of an endomorphism.

At a prime $\ell\ne p$, the same rational $p$ is a local
unit, fixes $1_{\mathbb Z_\ell}$, and preserves $V_\ell$.
Thus (CCF33) is localized at the original prime $p$. On the
crossed product its coefficientwise form is the defect of
$\operatorname{Ad}U_p$ on $B_p$, with group labels unchanged.
The real scaling $\vartheta_p$ is an endomorphism of every
sheaf and every quotient above, with the exact factor $p$
in (CCF27); the rational-prime action has the additional
boundary (CCF34). These two actions are connected by their
explicit adelic components, not identified by dropping that
boundary. The calculation supplies the full support and
mixed-place objects on which a global trace comparison must
now be performed.

\begin{figure}[ht]
\centering
\begin{tikzpicture}[node distance=16mm and 8mm,
 every node/.style={font=\footnotesize,align=center,text width=43mm},>=Stealth]
 \node (w) {$W$\\generic adelic coefficients};
 \node[right=of w] (d) {$D=\bigoplus_p W/W_p$\\all marked prime defects};
 \node[right=of d] (h) {$D/\delta W$\\mixed-place arithmetic image};
 \node[below=of d] (k) {$H_X^1(T_SX,U_S\mathcal H)$\\full split relative group};
 \node[below=of h] (s) {$W\otimes\widetilde H^0(S,\mathbb C)$\\additional support components};
 \draw[->] (w)--node[above] {$\delta$}(d);
 \draw[->] (d)--node[above] {quotient}(h);
 \draw[->] (d)--node[left] {canonical}(k);
 \draw[->] (k)--node[below] {onto}(s);
 \draw[->] (k)--node[above,sloped] {(CCF19)}(h);
\end{tikzpicture}
\caption{The actual Connes--Consani coefficient sheaf after the
split-base extension. The bottom row is completed on the left
by $0\to D$ in (CCF18); its kernel is exactly the marked defect
space. Every higher support cycle contributes
$W\otimes H^j(S,\mathbb C)$ in degree $j+1$. Equations
(CCF11)--(CCF19) prove the maps, and (CCF30)--(CCF32) compute
the attained positive arithmetic quotient. Human coefficient
source: Connes--Consani \cite{ConnesConsani2025}, Section 4.}
\end{figure}
